\subsection{Marco Conceptual}

\subsubsection{Computación en la nube y modelos de servicio}

``La computación en la nube es un modelo que permite el acceso ubicuo, conveniente y bajo demanda a un conjunto compartido de recursos de computación configurables, los cuales pueden ser rápidamente provisionados y liberados con una mínima gestión o interacción con el proveedor de servicios'' \cite{Mell2011}. Este modelo ofrece beneficios significativos como elasticidad, pago por consumo y alta disponibilidad, factores esenciales para la operación sostenible de plataformas como Paralegales.

La computación en nube se clasifica en distintos modelos de servicio, entre los cuales se destacan \textit{Software as a Service} (SaaS), \textit{Platform as a Service} (PaaS) e \textit{Infrastructure as a Service} (IaaS) \cite{Mell2011}. A estos se han sumado modelos más granulares como \textit{Backend as a Service} (BaaS), que ofrece a los desarrolladores funcionalidades listas para usar —bases de datos, autenticación, almacenamiento, notificaciones \textit{push}— sin necesidad de gestionar servidores, accesibles mediante un SDK; y \textit{Function as a Service} (FaaS), que permite ejecutar lógica de negocio en forma de funciones efímeras gestionadas por el proveedor. Firebase, empleado en Paralegales para autenticación y persistencia, es un ejemplo representativo de BaaS, mientras que AWS Lambda, utilizado para la capa de servicios del proyecto, ejemplifica el modelo FaaS. La coexistencia de ambos modelos en una misma plataforma se conoce como \textit{arquitectura híbrida} y constituye una de las decisiones de diseño centrales de este trabajo.

\subsubsection{AWS Well-Architected Framework}

El diseño de arquitecturas robustas y escalables en la nube se apoya en buenas prácticas consolidadas internacionalmente. AWS propone el \textit{Well-Architected Framework}, que articula seis pilares fundamentales para la construcción de sistemas seguros, eficientes y resilientes: excelencia operativa, seguridad, fiabilidad, eficiencia en el rendimiento, optimización de costos y sostenibilidad \cite{AWS2024}. Aunque este marco no impone una arquitectura específica, ofrece principios y preguntas guía que orientan decisiones técnicas en áreas críticas. En el contexto del presente proyecto, varios de sus pilares —especialmente seguridad, excelencia operativa y optimización de costos— informaron las decisiones de diseño tomadas, incluso cuando estas no involucraron la sustitución completa de un proveedor por otro, sino la introducción de capacidades complementarias sobre la arquitectura existente.

\subsubsection{Vendor lock-in y estrategias de mitigación}

En el contexto de la computación en nube, el término \textit{vendor lock-in} se refiere a la dificultad de cambiar de proveedor una vez que una organización ha adoptado profundamente su plataforma. Esta dependencia puede deberse a los costos de migración de datos, al uso de APIs propietarias o incompatibles con estándares abiertos, y a la integración estrecha con servicios específicos del proveedor. Según \textcite{OparaMartins2016}, esta dependencia puede traducirse en limitaciones de portabilidad, mayores costos a largo plazo y menor capacidad de adaptación a nuevas tecnologías. El trabajo de \textcite{Harauzek2022} amplía este análisis desde la perspectiva del usuario empresarial, identificando dimensiones técnicas, contractuales y económicas del fenómeno.

Frente a este riesgo, la literatura plantea estrategias de mitigación que no necesariamente exigen abandonar al proveedor incumbente. La introducción de capas de abstracción explícitas —por ejemplo, repositorios definidos por interfaces, capas de adaptación tipo \textit{Ports \& Adapters} o servicios intermediarios— permite encapsular las dependencias propietarias y preservar la opción de sustitución futura sin imponer una migración inmediata. En el caso de Paralegales, esta línea de mitigación se materializa mediante el patrón \textit{Handler--Service--Repository} en el \textit{backend} (que abstrae el motor de persistencia detrás de una interfaz) y la introducción de un \textit{Backend-for-Frontend} (que actúa como punto de control sobre las llamadas a servicios externos).

\subsubsection{Backend-for-Frontend (BFF)}

El patrón \textit{Backend-for-Frontend} fue formalizado por \textcite{NewmanBFF} y consiste en construir un \textit{backend} específico que sirve a un \textit{frontend} concreto, en lugar de exponer directamente un \textit{backend} de propósito general a todos los clientes. El BFF actúa como capa intermedia que adapta, agrega y compone llamadas a servicios subyacentes, ajustándolas a las necesidades particulares de la interfaz que consume sus respuestas. Entre sus ventajas se cuentan: la posibilidad de reubicar lógica sensible que de otro modo residiría en el cliente, la simplificación de la experiencia de desarrollo del \textit{frontend}, la centralización del manejo de errores y autenticación, y la habilitación de capacidades de observabilidad y auditoría sobre cada llamada saliente.

En la práctica contemporánea, frameworks como Nuxt~3 y Next.js facilitan la implementación de un BFF al ofrecer un entorno de ejecución del lado del servidor estrechamente integrado con la capa de presentación. Esto permite que ciertos flujos —especialmente los que requieren credenciales o lógica de orquestación cercana a la sesión del usuario— se ejecuten en el servidor sin abandonar el ecosistema de la aplicación cliente. En Paralegales, el BFF implementado sobre Nuxt~3 cumple esta función para flujos como la consulta a las APIs no documentadas de la Rama Judicial, donde actúa como \textit{proxy} controlado, evita restricciones CORS y centraliza el saneamiento, la auditoría y la observabilidad de las peticiones \cite{Newman2021}.

\subsubsection{Domain-Driven Design (DDD) y arquitectura modular}

El \textit{Domain-Driven Design} (DDD), propuesto por \textcite{Evans2003}, plantea construir el software en torno al dominio del negocio mediante el uso de un lenguaje ubicuo compartido entre técnicos y expertos del dominio, y mediante estructuras organizadas en \textit{bounded contexts} —fronteras explícitas dentro de las cuales un modelo del dominio mantiene consistencia y significado propio. Esta perspectiva combate la entropía típica de los sistemas que crecen sin un eje conceptual claro, donde la lógica de un dominio se mezcla con la de otros y los cambios se vuelven progresivamente más costosos.

Aplicado a una arquitectura \textit{frontend} moderna, DDD se traduce en una organización del código en módulos por dominio funcional, cada uno con sus componentes, servicios, modelos y rutas propios. En el ecosistema Nuxt, esta organización se materializa mediante \textit{Nuxt Layers}, un mecanismo que permite componer capas autónomas que actúan como módulos independientes y componibles. Cada capa puede tener su propia configuración, sus propias dependencias y su propio ciclo de evolución, sin acoplarse al resto del sistema. Esta estructura habilita la incorporación de nuevas \textit{features} sin riesgo de conflicto con el código existente y permite que cambios en un dominio no requieran análisis de impacto sobre el resto del proyecto.

\subsubsection{Arquitectura hexagonal y Repository Pattern}

La \textit{arquitectura hexagonal}, también conocida como \textit{Ports \& Adapters}, fue propuesta por \textcite{Cockburn2005} con el objetivo de aislar el núcleo de lógica de negocio de las preocupaciones de infraestructura (bases de datos, servicios externos, interfaces de usuario). En este modelo, la lógica de negocio define ``puertos'' —interfaces abstractas— que describen lo que el sistema necesita del exterior, mientras que ``adaptadores'' concretos implementan esos puertos contra tecnologías específicas. Este aislamiento permite que la lógica de negocio sea probada, modificada y razonada de forma independiente de la tecnología que la rodea.

El \textit{Repository Pattern}, descrito por \textcite{Fowler2002} y refinado en el marco de DDD por \textcite{Evans2003}, es la instancia más extendida del enfoque hexagonal aplicada a la persistencia. Consiste en definir una interfaz que abstrae las operaciones de almacenamiento (lectura, escritura, consulta) y delegar su implementación concreta a una clase específica para cada motor de persistencia. Esta abstracción aporta tres beneficios concretos: (i) habilita pruebas unitarias rápidas y deterministas mediante implementaciones \textit{mock}, sin acceder a una base de datos real; (ii) desacopla la lógica de negocio del proveedor de persistencia, preservando la opción de migrar el motor de datos sin tocar el resto del sistema; y (iii) clarifica las fronteras del dominio al exponer únicamente las operaciones que el negocio efectivamente requiere. En Paralegales, este patrón se aplica en la capa de servicios del \textit{backend} mediante la estructura \texttt{repository\_interface.go} + \texttt{repository\_firestore.go}.

\subsubsection{Arquitecturas \textit{event-driven} y mensajería asíncrona}

Los sistemas distribuidos modernos rara vez se comportan como cadenas síncronas de peticiones y respuestas. Cuando una operación implica trabajo costoso, integración con múltiples proveedores externos o procesamiento por lotes, el modelo síncrono degrada la experiencia del usuario y compromete la fiabilidad del sistema ante fallos parciales. La \textit{arquitectura event-driven}, basada en mensajería asíncrona, desacopla productores y consumidores mediante colas o tópicos: el productor publica un mensaje describiendo un evento y continúa su ejecución, mientras uno o varios consumidores procesan el mensaje de forma independiente.

El trabajo seminal de \textcite{Hohpe2003} cataloga los patrones de integración aplicables a este tipo de arquitecturas: colas punto a punto, tópicos de publicación/suscripción, enrutadores basados en contenido, transformadores de mensajes y filtros, entre otros. En el ecosistema AWS, estos patrones se materializan típicamente mediante Amazon SQS (colas) y Amazon SNS (tópicos), que en conjunto habilitan flujos asíncronos altamente desacoplados. En Paralegales, este modelo soporta casos de uso como el envío masivo de comunicaciones del módulo de campañas, donde una función orquestadora desglosa el volumen total de destinatarios y deposita instrucciones individuales en una cola, permitiendo que consumidores especializados procesen cada mensaje sin bloquear la experiencia del usuario administrador.

\subsubsection{Seguridad en aplicaciones web}

Garantizar la seguridad de las aplicaciones web es un desafío creciente en el entorno tecnológico actual. El \textit{OWASP Top 10} identifica las principales categorías de riesgo, entre las cuales se destaca \textit{Insecure Design}, una categoría introducida en la edición 2021 que se refiere a deficiencias estructurales en la arquitectura del software que permiten vulnerabilidades explotables \cite{OWASP2021}. A diferencia de las vulnerabilidades de implementación —que pueden resolverse parchando código defectuoso—, los problemas de \textit{Insecure Design} requieren rediseñar partes del sistema, lo que los hace particularmente costosos cuando se descubren tardíamente.

La exposición de lógica de negocio crítica en el \textit{frontend} —común en aplicaciones construidas con un enfoque ``cliente directo a BaaS''— constituye un ejemplo paradigmático de esta categoría: la lógica queda inspeccionable, manipulable y replicable por cualquier observador del cliente. \textcite{BirrEngwall2024} desarrollan el enfoque \textit{Secure by Design} como contrapartida, proponiendo el uso de \textit{domain primitives} y la reubicación sistemática de la lógica sensible en entornos de ejecución controlados. Esta línea de pensamiento es directamente aplicable a Paralegales, donde la introducción del BFF y la consolidación de la lógica sensible en funciones \textit{serverless} buscan precisamente reducir la superficie de exposición del cliente.

\subsubsection{Observabilidad de sistemas distribuidos y SRE}

En arquitecturas distribuidas, donde una sola acción del usuario puede atravesar múltiples servicios, redes y proveedores externos, el monitoreo tradicional resulta insuficiente. La \textit{observabilidad}, definida por \textcite{Sridharan2018}, es la capacidad de inferir el estado interno de un sistema a partir de sus salidas externas, y se articula clásicamente en torno a tres pilares: \textit{logs} estructurados, métricas agregadas y trazas distribuidas. La observabilidad no es una herramienta sino una propiedad del sistema, lograda mediante instrumentación deliberada, esquemas estandarizados y correlación de eventos a través de identificadores propagados.

La práctica de \textit{Site Reliability Engineering} (SRE), formalizada por Google \cite{Beyer2016}, complementa la observabilidad con un conjunto de principios operativos: definición explícita de \textit{Service Level Objectives} (SLOs), uso de \textit{error budgets} para balancear estabilidad y velocidad de cambio, y medición sistemática de indicadores como el \textit{Mean Time To Recovery} (MTTR). La reducción del MTTR depende directamente de la capacidad de detectar fallos rápidamente, diagnosticar su causa raíz y aplicar correcciones —tres capacidades que, a su vez, dependen de un sistema de observabilidad bien diseñado. En Paralegales, el subsistema de observabilidad implementado se inspira en estos principios: instrumentación mediante interceptores HTTP, esquema JSON estandarizado, métricas deducidas a partir de \textit{logs}, alertas proactivas y propagación de identificadores de correlación a través de la malla de servicios.

\subsubsection{Infraestructura como Código (IaC)}

La \textit{Infraestructura como Código} consiste en gestionar y aprovisionar recursos de infraestructura mediante archivos declarativos versionables, en lugar de mediante configuración manual a través de consolas web o comandos imperativos. \textcite{Morris2020} sostiene que IaC no es solo una técnica de automatización, sino un cambio de paradigma que permite tratar a la infraestructura con las mismas prácticas que el código de aplicación: revisión por pares, control de versiones, pruebas, despliegue continuo y reversión controlada. Esta perspectiva resulta especialmente valiosa en arquitecturas \textit{serverless}, donde la infraestructura activa cambia con frecuencia y cada función Lambda, cola SQS o alarma de CloudWatch debe estar versionada y reproducible.

En el ecosistema AWS, herramientas como CloudFormation, CDK y SAM permiten definir la infraestructura completa de un proyecto mediante archivos declarativos. En Paralegales, la totalidad de la infraestructura del repositorio \texttt{paralegales-services} —funciones Lambda, integraciones con API Gateway, colas SQS, tópicos SNS, dashboards de observabilidad y alertas— está definida mediante CloudFormation, lo que garantiza que cualquier entorno (desarrollo, \textit{staging}, producción) pueda reconstruirse de forma idéntica desde el repositorio.

\subsubsection{Prácticas SDLC para evolución continua}

La sostenibilidad técnica de un sistema a lo largo del tiempo depende menos de las decisiones aisladas tomadas en un momento dado y más de las prácticas que rigen su ciclo de vida (\textit{Software Development Life Cycle}, SDLC). \textcite{Martin2008} argumenta que la legibilidad y la claridad del código (\textit{Clean Code}) son determinantes para la mantenibilidad de un sistema, dado que los desarrolladores pasan considerablemente más tiempo leyendo código que escribiéndolo. Complementariamente, la comunidad de ingeniería de software ha consolidado un cuerpo de patrones de diseño orientados a objetos \cite{Gamma1994} y patrones de arquitectura empresarial \cite{Fowler2002} que ofrecen vocabulario común y soluciones reutilizables para problemas recurrentes.

A escala de equipo y de organización, las prácticas de \textit{Continuous Delivery} (CD), formalizadas por \textcite{Humble2010}, plantean automatizar todo el camino entre el \textit{commit} del desarrollador y el despliegue en producción, de modo que cada cambio pueda liberarse de forma rápida, repetible y de bajo riesgo. La CD se apoya en \textit{pipelines} de integración continua, baterías de pruebas automatizadas a múltiples niveles (unitarias, de integración, \textit{end-to-end}) y mecanismos de validación temprana como \textit{linters}, formateadores y \textit{hooks} de control de versiones. El concepto de \textbf{deuda técnica}, acuñado por \textcite{Cunningham1992}, describe la metáfora financiera asociada a tomar atajos durante el desarrollo: cada atajo es un préstamo que devenga ``intereses'' en forma de costo creciente para cada cambio futuro, hasta que la deuda se vuelve impagable. Las prácticas SDLC continuas son precisamente el mecanismo para evitar que la deuda crezca de forma descontrolada, manteniendo el sistema en condiciones de evolucionar sin ruptura.
